% =========================================================
% BAO CAO CHI TIET -- gui giao vien huong dan
% Bai 1: Du phong kha thi tinh nang moi trong he vi dich vu
%
% BIEN DICH:  pdflatex bao_cao_giao_vien.tex   (chay 2 lan cho muc luc)
% NEU LOI FONT TIENG VIET (thieu VnTeX): thay 2 dong inputenc/fontenc bang
%     \usepackage{fontspec}
%     \setmainfont{Times New Roman}
%   roi chay:  xelatex bao_cao_giao_vien.tex
%
% TODO TRUOC KHI GUI: dien ten + lop/khoa o \author
% =========================================================
\documentclass[11pt,a4paper]{article}
\usepackage[utf8]{inputenc}
\usepackage[T5]{fontenc}
\usepackage{amsmath}\usepackage{amssymb}\usepackage{booktabs}
\usepackage{array}\usepackage{geometry}\usepackage{hyperref}
\usepackage{xcolor}\usepackage{enumitem}
\geometry{margin=1in}
\hypersetup{colorlinks=true,linkcolor=blue,citecolor=red,urlcolor=teal}
\makeatletter\def\UTFviii@undefined@err#1{\relax}\makeatother
\newcommand{\luuy}[1]{\par\medskip\noindent\fcolorbox{gray!50}{gray!8}{%
  \parbox{\dimexpr\linewidth-2\fboxsep-2\fboxrule}{\small #1}}\par\medskip}
\newcommand{\sotrong}[1]{\textbf{#1}}

\title{\textbf{Dự phóng khả thi tính năng mới\\ trong hệ vi dịch vụ}\\[0.4em]
       \large Báo cáo chi tiết tiến độ nghiên cứu}
\author{[Tên sinh viên] --- [Lớp / Khoa]}
\date{07/10/2026}

\begin{document}
\maketitle

\begin{abstract}\noindent
Sau khi rà soát tài liệu rộng (kiến trúc hoá yêu cầu bằng LLM, dự báo hiệu năng
kiến trúc, suy luận nhân quả cho RCA, sinh workload bằng lý luận nhân quả, đặt vị trí
dịch vụ theo ý định ngôn ngữ), chúng tôi không tìm thấy công trình nào kết hợp đủ bốn
thuộc tính: nhận \textbf{ngôn ngữ tự nhiên}, dự báo \textbf{tiến cứu}, \textbf{không
cần} telemetry của chính tính năng đang được hỏi, và mang \textbf{bảo đảm cấu trúc}.
Một công trình tầm nhìn độc lập gần đây (arXiv 2605.02454, 2026) gọi tên đúng khoảng
trống này nhưng dừng ở ``chúng tôi đề xuất một benchmark agenda'' --- chưa triển
khai, chưa đo. Báo cáo này trình bày một lần triển khai và đo đạc tiến cứu đầu tiên
cho đúng bài toán đó: dự báo khả thi của một \emph{tính năng chưa được xây dựng}
trong hệ vi dịch vụ, bằng một giao thức đóng băng dự đoán + băm SHA-256 + agent cài
đặt bị cách ly khỏi mô hình, \emph{trước khi} đo đáp án thật. Dọc theo quá trình, nhóm tìm ra và \textbf{sửa xong, xác minh lại bằng đo đạc}
\sotrong{ba lỗi đo lường/phương pháp nghiêm trọng}: một lỗ hổng cho $8/11$
prompt đối kháng lọt qua lớp kiểm; một phát biểu kiểm được và sai; một lỗi âm
thầm biến lỗi hạ tầng thành ``an toàn giả'' trong benchmark. Trên nhánh tiến cứu
$n = 8$, sai số điểm gãy đạt \sotrong{12,2\%} (18,0 req/s), hơn hẳn hai mốc nền
trung thực. Báo cáo cũng nêu đầy đủ giới hạn của phương pháp hiện tại chưa giải
được bằng sửa mã: ba mệnh đề nghiên cứu bị bác bỏ, và một cạnh nhân quả (Tầng
2.5, backpressure) mà kẹp dải chỉ phục hồi được hành vi điển hình, còn để lại
rủi ro đuôi dài ở vài trường hợp cực trị --- bằng chứng đủ để, sau khi kiểm
thêm trên dữ liệu tải thật của chính nhóm, \textbf{gỡ hẳn cạnh này khỏi DAG sản
phẩm} thay vì tiếp tục vá. Đây là một giới hạn đặc trưng của một phép xấp xỉ
tuyến tính từng phần gần biên can thiệp, không phải một lỗi lập trình --- đúng
với tinh thần giao thức đánh giá mà báo cáo đề xuất: một đánh giá chỉ có giá
trị khi nó có khả năng tự bác bỏ.
\end{abstract}

\tableofcontents
\newpage

\section{Giới thiệu}
\textbf{Bài toán.} Cho một yêu cầu tính năng viết bằng ngôn ngữ tự nhiên và một tải
đỉnh kỳ vọng, trả lời: hệ vi dịch vụ \textbf{chưa có tính năng đó} còn đáp ứng SLO
không, và node nào sẽ nghẽn trước. Đây là \textbf{suy luận can thiệp} (what-if),
không phải dự báo chuỗi thời gian.

\textbf{Một bài toán chưa thấy ai đặt ra theo đúng tổ hợp này.} Rà soát tám họ công
trình liên quan gần nhất --- dịch yêu cầu NL sang kiến trúc bằng LLM; dự báo hiệu
năng kiến trúc (Palladio, LQN); suy luận nhân quả cho RCA (CIRCA, BARO); sinh
workload đối kháng bằng lý luận nhân quả; đặt vị trí dịch vụ theo ý định ngôn ngữ ---
không họ nào kết hợp đủ bốn thuộc tính: nhận \textbf{ngôn ngữ tự nhiên}, dự báo
\textbf{tiến cứu} (trước khi tính năng tồn tại), \textbf{không cần} telemetry của
chính tính năng đó, và mang \textbf{bảo đảm cấu trúc} cho đầu ra. Một công trình tầm
nhìn độc lập, gần đây (arXiv 2605.02454, 2026, ``Causal Software Engineering: A
Vision and Roadmap'') gọi tên đúng khoảng trống này --- coi việc triển khai phát
hành như một can thiệp $do(\cdot)$ và kêu gọi một ``benchmark agenda'' để đo tiến độ
--- nhưng tự nhận là tầm nhìn, chưa triển khai hay đo đạc gì (ngôn từ hiện tại của
bài: ``we \emph{outline}... we \emph{propose} three benchmark families''). Báo cáo
này là một lần triển khai và đo đạc \textbf{tiến cứu đầu tiên} cho đúng bài toán đó.

\textbf{Đóng góp phương pháp chính không nằm ở mô hình, mà ở thứ tự đo.} Dự đoán
được ghi ra tệp, băm SHA-256 và commit \emph{trước khi} tính năng được viết và
trước khi chạy đo. Người cài tính năng là một agent bị cách ly: không thấy taxonomy,
không thấy mô hình, không thấy kết quả vòng trước. Nếu phải sửa mô hình sau khi đã
thấy đáp án, vòng đo đó \textbf{bị tuyên bố là hồi cứu}, không được cứu.

\textbf{Cấu trúc báo cáo.} Mục~\ref{sec:tongquan} trình bày tổng quan hệ thống: hai
câu hỏi hệ trả lời, hợp đồng đầu vào, và các tiền điều kiện. Mục~\ref{sec:cautao} đi
vào cấu tạo từng phần: hai ngăn xếp Hệ A/Hệ B, đồ thị nhân quả ba tầng, và thang bóc
tách P0--P3 của Hệ B. Mục~\ref{sec:thucnghiem} trình bày các mệnh đề nghiên cứu, giao
thức đo tiến cứu, hạ tầng đo, và kết quả định lượng. Mục~\ref{sec:gh} nêu giới hạn và
các tuyên bố bị bác.

\section{Tổng quan hệ thống}\label{sec:tongquan}
Hệ trả lời \textbf{hai} câu hỏi, không phải một --- điểm dễ nhầm nhất khi đọc mã.

\begin{center}\small
\begin{tabular}{p{7.2cm}p{4.4cm}c}
\toprule Câu hỏi & Hàm & Cần thêm \\ \midrule
\emph{``Thêm tính năng này thì CPU/mem/latency từng service thành bao nhiêu?''}
 & \texttt{forecast\_resource\_impact} & --- \\
\emph{``Ở tải đỉnh $L$ req/s, hệ còn đáp ứng SLO không, gãy ở đâu?''}
 & \texttt{assess\_feasibility\_at\_peak} & $L_{\text{peak}}$ \\ \bottomrule
\end{tabular}\end{center}

Hai câu này là \textbf{hai tầng, không phải hai đối thủ}: phán quyết SLO bắt buộc
phải có dự báo tài nguyên trước.

\subsection{Hợp đồng đầu vào}
\begin{center}\small
\begin{tabular}{llp{6.2cm}}
\toprule Trường & Ai cung cấp & Ghi chú \\ \midrule
\texttt{requirement} & người dùng & mô tả ngôn ngữ tự nhiên \\
$L_{\text{peak}}$ & người dùng & \textbf{tường minh}; không lấy trung bình dữ liệu
  huấn luyện. Không cho thì hệ \textbf{từ chối kết luận} \\
$\lambda^{*}$ & mặc định + ghi đè & $=$ \texttt{anchor\%} $\times L_{\text{peak}}$;
  người dùng ghi đè được \\
$C_s$ (trần CPU) & hệ tự lấy & đọc trực tiếp từ Docker đang chạy \\
SLO & tệp cấu hình & p99 $\leq 250$ ms, lỗi $\leq 1\%$; \textbf{chốt ở Phase 0,
  không đổi sau khi đo đáp án} \\
$k_s$ & probe, tuỳ chọn & không truyền thì giả định $k_s = 1$ \\ \bottomrule
\end{tabular}\end{center}

\subsection{Ba tiền điều kiện: không thoả thì hệ từ chối, không đoán}
\begin{enumerate}[leftmargin=*,itemsep=2pt]
 \item \textbf{Phạm vi.} Yêu cầu phải hướng khách hàng và thuộc taxonomy; nếu không
   $\to$ \texttt{REFUSED} / \texttt{NEEDS\_HUMAN\_REVIEW}.
 \item \textbf{Hạn ngạch CPU.} Node nghẽn phải có $\geq 0{,}5$ core; nếu không $\to$
   \texttt{UNDECIDED}, vì dưới mức đó SLO vỡ do \emph{CFS throttling} chứ không do
   bão hoà CPU.
 \item \textbf{Khớp tên service} giữa đồ thị phụ thuộc và bảng số liệu.
\end{enumerate}

\luuy{\textbf{Cái bẫy nguy hiểm nhất là điều kiện thứ ba, vì nó không báo lỗi.} Bảng
số liệu gọi cổng vào một tên, đồ thị gọi tên khác $\Rightarrow$ cổng vào bị loại khỏi
đồ thị \emph{âm thầm}, mà hệ vẫn chạy trọn và vẫn in thông báo khớp thành công. Đã
thêm chốt cảnh báo; vẫn phải kiểm tay mỗi khi thêm hệ mới.}

\section{Cấu tạo từng phần}\label{sec:cautao}

\subsection{Hai ngăn xếp: Hệ A và Hệ B}
Luồng: NL $\to$ ParserAgent $\to$ (archetype, $\Delta$, chain) $\to$ \textbf{tầng dự
báo} ($W_s$, $\mathrm{CPU}_s$) $\to$ \textbf{tầng phán quyết} ($u_s$ so với $u^{*}$,
cổng SLO). Hàm thứ nhất dừng ở tầng dự báo; hàm thứ hai đi hết.

Hiện hai tầng bị \textbf{nhân bản trong hai ngăn xếp}, vì lý do \textbf{lịch sử}: hệ A
học từ dữ liệu fault-injection và đã sinh ra các số \emph{đã công bố}, nên đường khả
thi (hệ B) được viết tách ra để không chạm vào nó.

\paragraph{Hệ A.} DAG nhân quả 28 node (7 service $\times$ 4 metric), \sotrong{29 cạnh
cấu trúc} (Tầng 1 + Tầng 2 --- xem Mục~\ref{sec:tang25}). Cơ chế: hồi quy không âm
cho CPU/mem/workload; cơ chế hàng đợi trên $[X,\, X/(C-X)]$ với $C = P_{99}(X)\times
1{,}5$ cho latency. \textbf{Không còn cạnh Tầng 2.5}: bản trước có thêm 5 cạnh học
(tổng 34), nhưng đã gỡ hẳn khỏi \texttt{capacity\_agent.py} ngày 2026-10-10 sau các
phát hiện ở Mục~\ref{sec:tang25}; các số RQ1--RQ4 đã công bố \textbf{không đổi}
vì được sinh từ script đánh giá độc lập, không đọc lại DAG của \texttt{CapacityAgent}.

\paragraph{Hệ B.} Hai tầng \emph{phẳng}:
\[ \rho_s:\; W_s \sim W_{\text{gateway}} \text{ (qua gốc)}, \qquad
   \mathrm{CPU}_s = \alpha_s + \beta_s W_s, \quad \alpha_s,\beta_s \geq 0. \]
\textbf{Không có cạnh nào giữa các service} --- hình nan hoa, chain chỉ là \emph{mặt
nạ cộng tải}. \textbf{Không có node latency nào} --- đây là lý do cấu trúc khiến tiêu
chí khả thi buộc phải thuần CPU.

\subsubsection{Hai hệ đã được đo là một mô hình ở hai dạng}\label{sec:hainganxep}
Đo trên 456 hàng dữ liệu huấn luyện:
\begin{center}\small
\begin{tabular}{lc}
\toprule Đo ở mức & Lệch lớn nhất \\ \midrule
workload $W_s$, cả 7 service & 18,89\% \\
workload $W_s$, 5 node tính điểm & 13,25\% \\
\textbf{mức sử dụng $u_s$, 5 node tính điểm} & \sotrong{0,30\%} \\
\textbf{tại nút nghẽn (\texttt{front-end})} & \sotrong{0,00\%} \\ \bottomrule
\end{tabular}\end{center}

Service \texttt{user} lệch workload 13,25\% nhưng $\beta_{\text{user}} = 0{,}0674$ nhỏ
nên sai số đó gần như không dịch $u$. Nút nghẽn luôn là \texttt{front-end} --- gốc của
\emph{cả hai} cách tính --- nên lệch đúng 0,00\% và \textbf{phán quyết không đổi}. Vậy
chúng là \emph{một} mô hình viết ở hai dạng, \textbf{ở mức $u$}, là mức đi vào phán
quyết. Hợp nhất là \textbf{thiết kế mục tiêu} nhưng \textbf{chưa áp}: áp bây giờ là
đổi bộ dự đoán \emph{sau khi đã xem đáp án}, làm mọi số mới thành hồi cứu.

\luuy{\textbf{Tự phê.} Bản trước của đoạn này viết ``lệch $\leq 0{,}5\%$ trên mọi
service'' --- một phát biểu \emph{kiểm được và sai}, và không script nào tính ra nó.
Đã sửa và đã có script tái lập. Cùng mẫu lỗi với $R^2 = 0{,}0034$ ở Mục~\ref{sec:gh}.}

\subsubsection{Tận dụng cả hai hệ trong pipeline đang chạy, không hợp nhất mô hình}
Không hợp nhất Hệ A vào Hệ B (vì sẽ sửa bộ dự đoán sau khi đã xem đáp án, làm
hồi cứu mọi số đã đóng băng), nhưng 0,30\% ở trên là một con số \textbf{toàn
cục cho cả deployment} — thừa bảo thủ cho một yêu cầu cụ thể: yêu cầu nào
không chạm \texttt{user} thì không thể gặp đúng mức lệch đó. Đã cài một lớp
mỏng (\texttt{src/scm/doi\_chieu\_chain.py}, gọi từ
\texttt{feasibility\_agent.py}) đọc lại bảng đã tính sẵn và \textbf{lọc đúng
theo chain của từng yêu cầu}, loại \texttt{payment}/\texttt{shipping} khỏi
phép đối chiếu vì cùng lý do chúng không được Hệ B tính điểm (CPU quá nhỏ,
\texttt{shipping} xử lý bất đồng bộ — xem \texttt{DATA\_FRAMEWORK.md} mục 3).
\textbf{Không đổi `verdict`} — Hệ B vẫn quyết định số cuối; chỉ thêm một nhãn
độ tin đúng phạm vi.

Đo trên 8 tính năng làm vòng tiến cứu chính (bảng Mục~\ref{sec:ketquadinhluong}):
\begin{center}\small
\begin{tabular}{llc}
\toprule Tính năng & Chain đã tính điểm & Lệch lớn nhất (so với 0,30\% toàn cục) \\ \midrule
promo    & front-end, carts, orders                          & 0,078\% (chặt hơn) \\
recs     & front-end, user, catalogue, orders                & 0,301\% (bằng) \\
track    & front-end, orders                                 & 0,000\% (chặt hơn) \\
review   & front-end, user, catalogue                        & 0,301\% (bằng) \\
cartsum  & front-end, carts                                  & 0,078\% (chặt hơn) \\
quickadd & front-end, catalogue, carts                       & 0,078\% (chặt hơn) \\
express  & front-end, user, catalogue, carts, orders         & 0,301\% (bằng) \\
browse   & front-end, catalogue                              & 0,000\% (chặt hơn) \\ \bottomrule
\end{tabular}\end{center}

\textbf{5/8 tính năng nhận được một cận chặt hơn} con số toàn cục; 3 tính năng
còn lại (recs, review, express) đúng khi giữ nguyên 0,30\%, vì chain của
chúng chạm đúng \texttt{user} — nơi duy nhất còn phần dư. Đây \textbf{không
phải} cải thiện độ chính xác của phán quyết hay điểm gãy (hai số đó không
đổi) — là cải thiện \textbf{độ đúng phạm vi} của câu nói kèm theo phán quyết,
cho đúng yêu cầu đang được hỏi thay vì một cận xấu nhất mượn từ một service
yêu cầu đó chưa chắc chạm tới. Đã có test (\texttt{tests/test\_doi\_chieu\_chain.py})
chặn hồi quy nếu sau này Hệ A/Hệ B lệch xa hơn 1\%.

\subsection{Đồ thị nhân quả ba tầng: Tầng 2.5 được đề xuất, kiểm, rồi gỡ}
\begin{center}\small
\begin{tabular}{llp{5.4cm}}
\toprule & Cạnh & Nguồn gốc \\ \midrule
\textbf{Tầng 1} & $W_i \to W_j$ giữa service & \textbf{biết trước} từ sơ đồ gọi \\
\textbf{Tầng 2} & $W_j \to R_j$ cùng service & \textbf{biết trước} theo cấu trúc \\
\textbf{Tầng 2.5} & $R_i \to R_j$ & \textbf{phải HỌC từ dữ liệu} ---
  \emph{đã gỡ khỏi DAG sản phẩm, xem dưới} \\ \bottomrule
\end{tabular}\end{center}

Tầng 2.5 mô tả \textbf{backpressure đồng bộ}: callee nghẽn, thread pool của caller bị
giữ, caller nghẽn theo --- cơ chế \emph{không} nằm trong sơ đồ gọi. Phần dưới kể lại
toàn bộ quá trình đề xuất, kiểm, và cuối cùng \textbf{gỡ bỏ} cơ chế này, theo đúng thứ
tự nó đã xảy ra --- vì chính thứ tự đó là bài học phương pháp.

Chọn cạnh bằng 4 pha: liệt kê ứng viên từ chính đồ thị $\to$ chấm $R^2$ gain trên dữ
liệu giữ lại (67/33) $\to$ ngưỡng tự động bằng \emph{knee-point} $\to$ cổng an toàn
khi ngoại suy. Kết quả (trên RCAEval, dữ liệu dùng để huấn luyện Hệ A cho RQ1--RQ4):
\sotrong{16 ứng viên, 5 sống sót}.

\subsubsection{Ca cảnh báo --- phát hiện phương pháp quan trọng nhất của bài}\label{sec:tang25}
Trên dữ liệu giữ lại (đã là ngoài phân phối về tải), cạnh Tầng 2.5 cải thiện
\sotrong{23/23} node đủ điều kiện và giảm số ca đảo dấu từ 30 xuống 24. Trên bằng
chứng đó nó \emph{đã được triển khai}. Phép kiểm \textbf{can thiệp thật} ($n = 1080$)
bác bỏ nó:
\begin{center}\small
\begin{tabular}{lr}
\toprule Predictor & Suy giảm MAPE trung bình dưới can thiệp \\ \midrule
hai tầng $R_B \sim W_B$ & \sotrong{12,6} \\
Tầng 2.5 $R_B \sim R_A$ & \sotrong{1\,658,3} \quad ($\approx 132\times$ tệ hơn) \\
\bottomrule \end{tabular}\end{center}

\textbf{Nguyên nhân đã định vị:} cha của cơ chế là \emph{tài nguyên}, và dưới can
thiệp giá trị đó ra ngoài dải huấn luyện ở \sotrong{74,7\%} hàng, trong khi ở dữ liệu
giữ lại chỉ \sotrong{0,1\%}. Nghĩa là giao thức giữ lại \textbf{gần như không bao giờ
tạo ra điều kiện làm nó sai}. Trong dải, hai dạng \emph{không phân biệt được} (tỉ lệ
suy giảm trung vị $1{,}00\times$). \textbf{Cách sửa:} kẹp cha tài nguyên về biên dải
đã thấy lúc khớp mô hình; cha workload \emph{không} kẹp.

\luuy{\textbf{Một sai sót tinh vi đã gặp và sửa.} Bản đầu cài kiểu \emph{chuyển nhánh}
(``ra ngoài dải thì bỏ cha tài nguyên''). Nó \textbf{phá tính đơn điệu}: dự báo
\textbf{tụt 7,09 ngay tại biên dải}, vi phạm tiền điều kiện của mệnh đề bao chứng
nhận. Tệ hơn, phép kiểm tiền điều kiện \emph{không bắt được}, vì nó chỉ xét điều kiện
hệ số không âm --- mà hệ số của nhánh đầy vẫn không âm. Phải bổ sung phép kiểm đơn
điệu \emph{bằng số}. Sau khi đổi sang kẹp, tính trên cả $1\,080$ quan sát: \textbf{trung
vị} phục hồi ngang hai tầng (tỉ lệ $1{,}00$), nhưng \textbf{trung bình} còn
\sotrong{30,6} --- cao hơn hai tầng (12,6) khoảng $2{,}4\times$, vẫn tốt hơn bản
không kẹp (1\,658,3) tới $54\times$. Phần dư tập trung ở vài cặp (service, loại
lỗi) cụ thể --- nặng nhất \texttt{user}$\to$\texttt{front-end} CPU dưới packet
loss (4\,580 điểm MAPE). Đây là \textbf{giới hạn đặc trưng của một phép xấp xỉ
tuyến tính từng phần gần biên can thiệp cực trị}: kẹp khôi phục hành vi điển
hình về ngang mức nền, nhưng không loại bỏ hoàn toàn rủi ro đuôi dài ở vài
trường hợp hiếm.}

\paragraph{Quyết định cuối: gỡ hẳn, không chỉ kẹp.} Kẹp dải là một cải thiện bền,
nhưng để lại đúng rủi ro đuôi dài ở trên. Thay vì tiếp tục vá, nhóm chạy lại
\emph{đúng} giao thức chọn cạnh bốn pha (R$^2$-gain + knee-point + OOD-safety) --- lần
này trên dữ liệu tải \textbf{thật do nhóm tự đo} (\texttt{SS-TRAIN}), không phải
RCAEval. Kết quả: \sotrong{0/8} ứng viên Tầng 2.5 sống qua nổi vòng lọc đầu tiên.
Lý do định vị được: trên \texttt{SS-TRAIN}, hai tầng \emph{đã} giải thích được
$0{,}74$--$0{,}9972$ phương sai của mỗi node tài nguyên, so với chỉ $\approx 0{,}015$
trên RCAEval --- không còn phần phương sai nào để Tầng 2.5 "giải thích thêm" một
cách có ý nghĩa, kể cả giả tạo. Hai bằng chứng (can thiệp thật bác bỏ Tầng 2.5 thô;
chọn cạnh trên dữ liệu thật của chính hệ cũng không chọn nó) cùng chỉ về một hướng,
nên Tầng 2.5 đã \textbf{bị gỡ hẳn khỏi \texttt{capacity\_agent.py}} (2026-10-10),
không giữ lại dưới dạng kẹp nữa. \texttt{src/scm/co\_che\_cong\_dai.py} (cơ chế kẹp)
vẫn còn trong repo làm tài liệu lịch sử, nhưng không còn script nào gọi nó từ DAG sản
phẩm. \textbf{Không ảnh hưởng} các số RQ1--RQ4 đã công bố: chúng được sinh từ
\texttt{evaluation\_suite.py} và các script RQ độc lập, vốn tự xây cạnh backpressure
riêng, không đọc lại DAG của \texttt{CapacityAgent}.

\paragraph{Bài học tổng quát.} \emph{Độ chính xác trên dữ liệu giữ lại dưới một giao
thức quan sát không dự đoán được hành vi dưới can thiệp, kể cả khi giao thức giữ lại
đó đã là ngoài phân phối.} Thuộc tính phân biệt hai tầng với Tầng 2.5 là \emph{cha
workload ra ngoài dải ít hơn hẳn cha tài nguyên} --- một \emph{prior về độ bền}, không
phải một lập luận nhận dạng nhân quả. Bài học thứ hai, từ chính việc gỡ: \textbf{một
cơ chế có thể đúng trên một nguồn dữ liệu (RCAEval) và vô nghĩa trên nguồn khác
(chính hệ đang vận hành)}, vì hai tầng đã chiếm hết phương sai cần giải thích --- nên
``giữ cạnh vì nó cải thiện ở đâu đó'' không phải lý do đủ để giữ nó trong DAG sản
phẩm của \emph{hệ này}.

\subsection{Thang bóc tách P0 đến P3 (Hệ B)}
Các $P$ \textbf{không phải mô hình cạnh tranh} --- chúng là \textbf{thang bóc tách},
mỗi bậc thêm \emph{đúng một} giả định, nên hiệu số giữa hai bậc quy được sai số về một
nguyên nhân.
\begin{center}\small
\begin{tabular}{llll}
\toprule & Công thức tải & Thêm gì & Tư cách \\ \midrule
P0 & $W_s = \rho_s L (1+\Delta)$ & --- (rải đều) & tiền đăng ký \\
P1 & $W_s = L(\rho_s + \Delta k_s [s \in \text{chain}])$ & chain làm mặt nạ & tiền đăng ký \\
P1\_ctrl & như P1, chain \emph{ngẫu nhiên} & đối chứng & tiền đăng ký \\
P2 & thêm chi phí mỗi lượt gọi & chi phí riêng từng service & tiền đăng ký \\
P3 & $=$ P2 với $k_s$ \emph{đo bằng probe} & bội số gọi thật & hồi cứu \\ \bottomrule
\end{tabular}\end{center}

Phán quyết: $u_s = \mathrm{CPU}_s / C_s$ trên 5 node, ngưỡng $u^{*} = 0{,}8783$ hiệu
chuẩn từ điểm gãy ramp nền. $R^{*}$ có nghiệm đóng vì mọi thứ tuyến tính theo $L$.

\subsubsection{Kết quả: nội dung nhân quả nằm ở trọng số, không ở topology}
Sai số mức sử dụng node trong chain (điểm \% của trần):
\begin{center}\small
\begin{tabular}{lrl}
\toprule Biến thể & Sai số & Diễn giải \\ \midrule
P0 rải đều & 7,49 & \\
P1 chain thật & 6,46 & $\leftarrow$ topology mua được \sotrong{1,03} \\
P1 chain SAI & 7,02 & \\
P3 chain $+\,k$ & 1,46 & $\leftarrow$ \sotrong{TRỌNG SỐ} mua được \sotrong{5,00} \\
\bottomrule \end{tabular}\end{center}

Chain \emph{thật} so với chain \emph{sai}: \textbf{không tách được} (Wilcoxon $p$
trung vị $0{,}19$ qua 200 lần bốc, \sotrong{8/16} ô hoà). Lý do rất cứng: \textbf{nút
nghẽn luôn là gateway, mà gateway nằm trong \emph{mọi} chain theo định nghĩa.}

\section{Các mệnh đề nghiên cứu, thực nghiệm, và phương pháp thực nghiệm}\label{sec:thucnghiem}
Đây không phải các mục cần giấu --- chúng là bằng chứng cụ thể nhất cho việc giao
thức đánh giá trong báo cáo này \emph{hoạt động đúng như thiết kế}: phát hiện sai số
trước khi nó lọt vào số liệu công bố.

\subsection{Ba lỗi nghiêm trọng --- tìm ra, sửa, và xác minh lại bằng đo đạc}
\begin{enumerate}[leftmargin=*,itemsep=3pt]
  \item \textbf{Lỗ hổng cổng kiểm phạm vi.} $8/11$ prompt đối kháng khớp đủ từ khoá
    với một archetype \emph{sai} vẫn lọt qua lớp kiểm. \textbf{Sửa:} gỡ hẳn đường tắt
    bỏ qua LLM. \textbf{Xác minh:} tỉ lệ lọt về $0/11$ qua $61$ lần gọi thật.
  \item \textbf{Một phát biểu kiểm được và sai.} Bản trước ghi ``hai ngăn xếp lệch
    $\leq 0{,}5\%$'' --- không script nào tính ra số đó. \textbf{Sửa:} đo lại, viết
    script tái lập. \textbf{Kết quả đúng:} $0{,}30\%$ / $0{,}00\%$ tại nút nghẽn
    (Mục~\ref{sec:hainganxep}).
  \item \textbf{Lỗi âm thầm biến thất bại hạ tầng thành ``an toàn giả''.} Một
    \texttt{except Exception} quá rộng trong benchmark parser bắt luôn lỗi 429 hết
    hạn ngạch, trả về đúng bộ ba giá trị tình cờ an toàn theo cả ba chỉ số đang đo
    --- khiến $96$--$98\%$ số dòng của một lượt đo bị đọc nhầm là ``LLM an toàn
    tuyệt đối''. \textbf{Sửa:} tách lời gọi LLM ra ngoài \texttt{try}, chỉ bắt lỗi
    parse JSON. \textbf{Xác minh:} chạy lại gặp \emph{đúng} hạn ngạch cũ, lần này
    dừng sạch và báo rõ nguyên nhân thay vì âm thầm nhiễm dữ liệu.
\end{enumerate}

\subsection{Trạng thái các mệnh đề nghiên cứu}
Áp cùng kỷ luật đó lên 5 mệnh đề đã công
bố trước đó: \sotrong{2 sống, 3 bị bác} --- ba trong số các phản bác nhắm vào kết quả
\emph{chính nhóm đã công bố trước đó}. Khác với ba lỗi ở trên, \textbf{ba phản bác
này không phải lỗi để sửa bằng mã} --- chúng là biên thật của phương pháp hiện tại,
đo được và không đổi khi đo lại. Một đánh giá chỉ có giá trị khi nó có khả năng tự
bác bỏ; vì dự đoán đã niêm phong trước, mỗi phản bác cũng là bằng chứng tiến cứu về
giới hạn của mô hình, không phải một lần đo hỏng.

\begin{center}\small
\begin{tabular}{clp{7.2cm}}
\toprule & Mệnh đề & Kết cục \\ \midrule
1 & Necessity & \textbf{bác} --- sai số biên độ giảm 25$\times$ nhưng chỉ sai cổng
    vào vẫn 92\% \\
2 & Guarantee & \textbf{bác cả hai} nhánh \\
3 & Measurability & \textbf{bác} --- benchmark \emph{đảo ngược} thứ hạng mô hình \\
4 & Attribution & đồ thị \textbf{sống}; toán tử $do(\cdot)$ \textbf{bác} \\
5 & Prospective & một phần sống, một phần bác \\ \bottomrule
\end{tabular}\end{center}

\subsection{Giao thức đo tiến cứu}
Thứ tự bất khả xâm phạm: (1) agent mù cài tính năng; (2) probe đo $k_s$, chi phí
đọc/ghi, độ trễ khi rảnh --- \textbf{không dùng dữ liệu tải}; (3) \textbf{đóng băng +
hash} dự đoán, commit git trước khi đo; (4) \emph{chỉ sau khi có hash} mới chạy ramp;
(5) đối chiếu.

\subsubsection{Cách ly agent cài đặt}
Agent \textbf{chỉ thấy}: mã nguồn gốc front-end, danh mục API backend \emph{trung lập}
(không nhóm theo archetype), yêu cầu ngôn ngữ tự nhiên kèm giao diện HTTP bắt buộc.
Agent \textbf{không thấy}: taxonomy chuỗi gọi, các tệp đã đóng băng, tài liệu hệ
thống, bản nháp bài, kết quả vòng trước, \textbf{và mã của các tính năng đã cài trước
đó}. Cài trong sandbox ngoài repo; agent \textbf{không được chạy hệ thống}, chỉ viết
mã; lỗi chức năng chuyển lại cho agent, \emph{không} chuyển gợi ý thiết kế. Mã sandbox
được băm SHA-256.

\luuy{\textbf{Quyết định thiết kế thực nghiệm đáng nêu.} 6 tính năng $\times$ 1 cường
độ \emph{hơn} 3 tính năng $\times$ 2 cường độ, dù cùng số ô. Hai cường độ của
\emph{cùng} một tính năng dùng chung $k_s$, chung chain, chung bản cài --- không phải
hai quan sát độc lập. Người phản biện sẽ gọi đó là pseudo-replication và tính lại $n$.}

\subsubsection{Các chốt từ chối đã cài trong mã}
Đây là chốt trong mã nguồn, không phải quy ước trên giấy: thiếu ánh xạ archetype $\to$
script đóng băng \textbf{từ chối chạy}; thư mục ramp đã có dữ liệu của tính năng đó
$\to$ \textbf{từ chối} (đã nhìn trước đáp án); $u^{*}$ chỉ hiệu chuẩn từ ramp nền; cơ
chế phải được kiểm trên các mức tải \emph{giữ lại} trước khi dùng để dự đoán.

\subsection{Hạ tầng đo và định nghĩa điểm gãy}
14 container, \textbf{không} có bộ mô phỏng người dùng. Bộ thu chạy \emph{cùng mạng
Docker}, lấy CPU từ bộ đếm cgroup và workload/latency từ endpoint số liệu \emph{thật}
của từng service; xuất 4 dòng ra đúng khuôn bộ dữ liệu chuẩn. Bộ sinh tải \textbf{mở
(Poisson), chạy trên HOST} --- để CPU của chính nó không lẫn vào CPU đang được đo.

\luuy{\textbf{Hai hướng đã thử và bỏ.} (i) Đưa bộ sinh tải vào container gắn thẳng vào
mạng Docker: cùng mức tải, bản chạy ngoài đạt SLO còn bản trong \emph{vỡ nặng} --- lỗi
khi đó là của harness, không phải của hệ. (ii) Dùng \texttt{docker stats} và đếm dòng
log làm mức tải: \emph{sai} với service Java/Go vì chúng không ghi log từng request.}

\textbf{Lưới tải hình học} $40, 50, 65, 80, 100, 125, 155, 195, 240$ req/s (bậc nhân
khoảng 1,25). Lưới cộng 20 req/s cho độ phân giải tương đối 10\% ở điểm gãy 200 nhưng
\sotrong{50\%} ở điểm gãy 40, nên sai số tương đối \textbf{không so sánh được giữa các
ô}. Mỗi bậc giữ tải một khoảng, \textbf{bỏ 40\% đầu} làm khởi động nóng, rồi chấm
$p_{99}$ và tỉ lệ lỗi trên phần ổn định.

\textbf{Điểm gãy} là \textbf{bậc đầu của chuỗi vi phạm liên tiếp cuối cùng}. Một bậc
vi phạm rồi hồi lại (dọn rác, nhiễu) ghi riêng là \emph{transient}, \textbf{không}
phải điểm gãy. Kết quả là một \textbf{khoảng}, không phải một số.

\textbf{Tiền đăng ký ở cấp phép đo.} Manifest ghi \textbf{trước khi đo}: SLO, cấu hình
trần, danh sách bậc tải, kế hoạch đầy đủ, mọi tham số, và \textbf{SHA-256 của chính
script đo}. Thứ tự ramp bị \emph{xáo ngẫu nhiên} để mức tải không lẫn với trôi theo
thời gian.

\subsubsection{Ba cái bẫy đã thực sự xảy ra}
\begin{enumerate}[leftmargin=*,itemsep=3pt]
 \item \textbf{Nhiễm trạng thái giữa các ramp.} Pool tài khoản dùng chung, không đặt
   lại. Tính năng \emph{sinh} đơn làm tính năng \emph{đọc} lịch sử đơn chậm dần theo
   thứ tự chạy: $155\text{--}195 \to 125\text{--}155 \to 80\text{--}100 \to$
   \textbf{vỡ ngay ở 40 req/s}. Điểm gãy đo được khi đó là của \emph{thứ tự chạy},
   không của tính năng. Đã chặn bằng hàm \textbf{từ chối chạy}.
 \item \textbf{Cờ lặp lại không có tác dụng} ở chế độ ramp, và mỗi lần chạy đánh số
   lượt lại từ 1 nên chạy vào thư mục cũ sẽ \textbf{ghi đè} dữ liệu.
 \item \textbf{Volume cơ sở dữ liệu tích luỹ.} Sau một đêm: 134\,871 đơn, 79\,836
   khách, 8\,314 giỏ, sống sót qua mọi lần khởi động lại $\Rightarrow$ điểm gãy nền
   \textbf{tụt từ 240 còn 125} (ổn định 4/4 ramp, không phải nhiễu).
\end{enumerate}
Hậu quả thật: một chiến dịch 18 ramp bị \textbf{loại hoàn toàn} vì hệ đã trôi trong
đêm (điểm gãy nền $195 \to 155 \to 155 \to 125$).

\textbf{Hợp đồng dữ liệu được cưỡng chế bằng mã}, thoát mã khác 0 nếu FAIL: cột bắt
buộc; đơn vị CPU hợp lý (bắt lỗi 100$\times$); tập huấn luyện \emph{không} được có
dòng tính năng / trần / khối bão hoà; mỗi ramp phải có điểm gãy và manifest tiền đăng
ký; liệt kê các ramp thuộc \textbf{tập khoá} không được dùng khi chỉnh mô hình.

\subsection{Kết quả định lượng}\label{sec:ketquadinhluong}

\subsubsection{Lớp đọc yêu cầu (parser)}
$\Delta$ \textbf{không} do LLM ước lượng tự do: $\Delta = \text{neo(archetype)} +
\text{điều chỉnh}$, với ba chốt theo thứ tự --- độ khớp thấp $\Rightarrow$ điều chỉnh
$:= 0$; $|\text{điều chỉnh}| > 10$ điểm $\Rightarrow$ \textbf{lùi hẳn về neo}; cuối
cùng kẹp $\Delta \in [5,50]$. Vì vậy $\Delta \in [5,50]$ \emph{bất kể} LLM trả gì.

Benchmark 50 prompt, 3 lần chạy LLM thật (Anchor MAE bằng \emph{điểm phần trăm}):
\begin{center}\footnotesize
\begin{tabular}{lcccc}
\toprule Cấu hình & Vi phạm biên & Bịa service & Sai cổng vào & Anchor MAE \\ \midrule
LLM trần, zero-shot & 38,0$\pm$2,0\% & 91,3$\pm$1,2\% & 100,0$\pm$0,0\% & 844,6$\pm$568,2 \\
few-shot, có bảng hiệu chuẩn, không chốt & 14,7$\pm$1,2\% & 2,0$\pm$0,0\% & 84,7$\pm$1,2\% & 38,6$\pm$0,2 \\
chỉ luật, không LLM & 0,0\% & 0,0\% & 0,0\% & 2,8$\pm$0,0 \\
\textbf{có chốt kiểm (của bài)} & \sotrong{0,0\%} & \sotrong{0,0\%} & \sotrong{0,0\%}
 & \sotrong{1,6$\pm$0,1} \\ \bottomrule
\end{tabular}\end{center}

Lặp trên backend thứ hai (117B, trọng số mở) tách được hai thứ mà một backend gộp lại.
\textbf{Năng lực model mua được độ hợp lý ngữ nghĩa}: Anchor MAE của LLM trần giảm
\sotrong{25$\times$} ($844{,}6 \to 33{,}6$). \textbf{Không mua được tính đúng cấu
trúc}: sai cổng vào vẫn \sotrong{92,0\%}, bịa service vẫn \sotrong{86,0\%}, và không
tập trung ở nhóm đối kháng (90\%, 100\%, 100\%, 80\% theo nhóm) --- vi phạm ràng buộc
topo một cách \emph{hệ thống}, không do prompt khó.

\luuy{\textbf{Hai điều phải nói kèm.} (a) Ba cột $0\%$ là \emph{do cấu trúc}: chốt
kiểm là hàm toàn phần nên tỉ lệ vi phạm bằng 0 theo định nghĩa. Vì vậy phép so với
baseline \emph{không} chứng minh chốt kiểm gánh việc; nó chỉ xác lập tiền đề rằng một
generator không bị chặn vi phạm những ràng buộc đó ở tỉ lệ cao \emph{dù được cho đúng
cùng ngữ cảnh}. Do đó bài báo \textbf{tỉ lệ kích hoạt} thay thế: 18/18 lần sửa cổng
vào, 33/36 lần kẹp biên. (b) Con số $1{,}6$ tính trên \sotrong{121/150} lời gọi: cổng
phạm vi trả 18 ($12{,}0\%$) cần người xem lại, 11 ($7{,}3\%$) từ chối.}

\paragraph{Một lỗ hổng đã tìm ra và sửa.} Đường tắt bỏ qua LLM khi độ khớp từ khoá cao
đã bị \textbf{gỡ hẳn}, sau khi phát hiện \sotrong{8/11} prompt đối kháng khớp đủ từ
khoá với \emph{một archetype sai} vẫn kích hoạt được đường tắt --- và đường tắt gán
cứng cờ tự khai báo thành ``đúng'' vô điều kiện, nên lớp kiểm phạm vi \textbf{không
bao giờ được thực thi} cho chúng. Tỉ lệ lọt thật là $72{,}7\%$, không phải $0\%$ như
kiểm chứng trước --- kiểm chứng đó gọi thẳng hàm con, không qua đường tắt.

\subsubsection{Sai số điểm gãy --- nhánh tiến cứu, \texorpdfstring{$n = 8$}{n = 8}}
\begin{center}\small
\begin{tabular}{lrrrr}
\toprule Mô hình & Sai số TB & Trung vị & Xấu nhất & TB (req/s) \\ \midrule
P1 --- mốc tiền đăng ký & 19,4\% & 20,3\% & 34,1\% & 28,0 \\
\textbf{P2} & \sotrong{12,2\%} & 6,3\% & 29,9\% & \sotrong{18,0} \\
P3 --- $k$ đo được & 12,6\% & 8,8\% & 29,9\% & 18,3 \\
\texttt{const\_base} --- bỏ qua tính năng & 27,4\% & 32,5\% & 52,9\% & 39,1 \\
\texttt{const\_oracle} --- hằng số tốt nhất & 11,5\% & 10,5\% & 24,4\% & 19,1 \\
\bottomrule \end{tabular}\end{center}

Wilcoxon ghép cặp (sàn kiểm định ở $n=8$ là $p = 0{,}0078$): \textbf{P2 hơn P1} và
\textbf{P2/P3 hơn \texttt{const\_base}} ($p = 0{,}0156$) $\Rightarrow$ tuyên bố được.
\textbf{P2/P3 KHÔNG hơn \texttt{const\_oracle}} ($p = 1{,}000$). \textbf{P3 không hơn
P2} ($p = 0{,}875$) --- hiệu ứng bội số $k$, đóng góp chính của vòng trước,
\emph{không lặp lại}.

\begin{center}\small
\begin{tabular}{lcccc}
\toprule Tập & $n$ & P1 & P2 & P3 \\ \midrule
dev & 4 & 37,6\% & 14,4\% & 15,3\% \\
\textbf{khoá} & 4 & 9,6\% [5,0--14,2] & \sotrong{2,7\% [2,2--3,3]} & 2,7\% \\
độc lập & 5 & 119,7\% & 87,5\% & 72,6\% [15,9--151,9] \\
tiến cứu vòng 1 & 2 & 10,0\% & 14,0\% & --- \\ \bottomrule
\end{tabular}\end{center}

Tập \textbf{khoá} chắc chắn nhất. Tập \textbf{độc lập} có khoảng \emph{cực rộng} ---
con số $72{,}6\%$ \textbf{không nên trích dẫn đơn lẻ}.

\paragraph{Hướng nguy hiểm (báo riêng).} Số ô mà mô hình dự đoán hệ còn ``sống'' ở
đúng mức tải đã làm vỡ SLO: \textbf{P1 3/8} $\to$ \textbf{P2 2/8} $\to$ \textbf{P3
2/8}. Phải báo tách khỏi sai số trung bình: một mô hình sai số trung bình đẹp vẫn có
thể báo KHẢ THI ở đúng mức tải đã làm vỡ hệ --- lỗi tốn kém nhất.

\subsubsection{Bài 2 --- định vị nguyên nhân (mục ngắn)}
Mệnh đề: độ chính xác định vị do \textbf{tầng metric được soi} quyết định, không do
thuật toán xếp hạng. Dữ liệu: 270 ca tiêm lỗi thật, 3 hệ thống. Một hệ, nhóm lỗi
\emph{mạng}, tập ứng viên mức ứng dụng:
\begin{center}\small
\begin{tabular}{lccc}
\toprule Tầng soi & Thuật toán A & Thuật toán B & Sàn ngẫu nhiên \\ \midrule
chỉ \texttt{cpu} & 13,3\% & 26,7\% & 9,1\% \\
chỉ \texttt{latency} & 60,0\% & 70,0\% & 10,0\% \\
mọi cột & 76,7\% & 80,0\% & 9,1\% \\ \bottomrule
\end{tabular}\end{center}
Đổi \emph{thuật toán} được $+3$ đến $+13$ điểm; đổi \emph{tầng} được \sotrong{$+47$
điểm} với cùng thuật toán. Với nhóm lỗi \emph{tài nguyên} thì \texttt{cpu} lại là tầng
đúng (83,3\% và 96,7\%). Đối chứng âm (không đọc dữ liệu) nằm \textbf{đúng trên sàn}.

\luuy{\textbf{Cần kiểm lại trước khi đưa vào bài.} Hai con số hay được nhắc --- 91,5\%
độ chính xác hạng nhất, và một baseline đã xuất bản được nâng từ 24,4\% lên 85,9\% ---
hiện \emph{chỉ có trong thông điệp commit}, và hai tài liệu của bài 2 \emph{tự đánh
dấu là lỗi thời}. Phải tính lại từ dữ liệu kết quả.}

\section{Giới hạn và các tuyên bố bị bác}\label{sec:gh}
\begin{enumerate}[leftmargin=*,itemsep=5pt]
\item \textbf{Không hơn một hằng số chọn khéo.} Dải điểm gãy thật của 8 tính năng chỉ
từ \sotrong{127,5 đến 195 req/s} --- đủ hẹp để một hằng số ở giữa cũng ngang ngửa. Dữ
liệu \emph{chưa} chứng minh mô hình \textbf{phân biệt được giữa các tính năng}, mới
chứng minh nó \textbf{biết mức điển hình}. \texttt{const\_oracle} dùng chính đáp án
nên không dùng được thực tế; nó là \emph{cận dưới} của mọi hằng số.

\item \textbf{Sàn nhiễu khoảng 13\%, ngang với chính sai số.} Biên độ giữa hai lần lặp
cùng một ô là \sotrong{21,4\%} (trung vị); ô tệ nhất lệch \sotrong{43\%}. Quy ra sai
số chuẩn mỗi ô khoảng 13\%, \emph{cùng bậc} với sai số 12,6\%. Hệ quả: sai số
\emph{tuyệt đối} \textbf{không phân giải được} dưới ngưỡng đó; chỉ so sánh \emph{ghép
cặp} còn hợp lệ, vì nhiễu là chung và bị khử khi lấy hiệu.

\item \textbf{Sàn phân giải --- giới hạn không thể vượt.} Hai tính năng khai báo
\emph{y hệt nhau} (cùng chuỗi gọi, cùng neo 15\%) có năng lực đo được \sotrong{70 so
với 125 req/s}. Cùng đầu vào thì \emph{buộc} cùng đầu ra, nên mô hình \textbf{sai ít
nhất 1/3} ở một trong hai, bất kể chỉnh tham số thế nào. Đây là tính chất của
\emph{thông tin có sẵn}, không phải của chất lượng mô hình. Đo chuỗi gọi thật
\emph{không} cứu được --- nó lệch \textbf{ngược chiều}. Khác biệt thật nằm ở \emph{giá
mỗi lời gọi}, không ở \emph{số} lời gọi.

\item \textbf{Neo $\Delta$ là số viết tay, và nằm ở cả hai phía của phép so sánh.}
Phía dự đoán dùng giá trị neo của archetype; phía đo bơm traffic tính năng ở \emph{đúng
cùng con số đó}. Nên vòng đo tiến cứu --- dù niêm phong, dù băm, dù agent mù ---
\textbf{chưa bao giờ kiểm chính cái neo}. Tuyên bố đúng là \emph{``cho trước $\Delta$,
hệ định vị được điểm gãy''}, hẹp hơn \emph{``hệ dự phóng được tác động của một tính
năng mới''}. Bảng phân chia trách nhiệm của bài xếp đại lượng này là \emph{không tự
động hoá được}. \textbf{Cần viết vào mục Threats to Validity --- hiện chưa có.}

\item \textbf{Throttling nằm ngoài phạm vi.} Tiền đề ``SLO vỡ ở mức sử dụng khoảng
0,88'' \emph{chỉ} kiểm chứng cho hạn ngạch $\geq 0{,}5$ core. Hai cấu hình trần nhỏ
đều thất bại: SLO vỡ ở mức sử dụng trung bình $20\text{--}45\%$ do bị bóp theo chu kỳ.
Khoảng $(0{,}15;\,0{,}5)$ core \textbf{chưa đo}.

\item \textbf{$n$ nhỏ ở mọi nhánh} (tiến cứu vòng 1 $n=2$, vòng 2 $n=8$ đến 10; khoá
$n=4$; độc lập $n=5$; 7 service). \textbf{Một hệ, một máy, một topology}: gateway luôn
nghẽn nên \emph{không kiểm được} câu ``node nào nghẽn''.

\item \textbf{Một con số trong Abstract chưa tái lập được.} $R^2 = 0{,}0034$ là
\emph{hằng số viết cứng}, không script nào tính ra nó. Tính lại cho $0{,}0150$ /
$0{,}0010$ / $0{,}0204$. \textbf{Kết luận không đổi} (mọi giá trị $\leq 0{,}02$, nhỏ
hơn hẳn $0{,}3009$) nhưng \textbf{con số phải được làm lại trước khi nộp}.
\end{enumerate}

\section{Việc còn lại trước khi nộp}
Xếp theo mức đe doạ việc được nhận:
\begin{enumerate}[leftmargin=*,itemsep=2pt]
 \item \textbf{Làm lại $R^2 = 0{,}0034$} cho có nguồn gốc --- số này nằm trong Abstract.
 \item \textbf{Tính lại các số của bài 2} từ dữ liệu kết quả; cập nhật hai tài liệu
   đang lỗi thời.
 \item \textbf{Viết neo $\Delta$ thành ranh giới phạm vi tường minh} trong Threats to
   Validity. Nói trước làm bài \emph{mạnh lên}, không yếu đi.
 \item \textbf{Viết kết quả tương đương hai ngăn xếp} ($0{,}30\%$ / $0{,}00\%$) vào
   bài; hợp nhất mã để làm \emph{future work}.
 \item \textbf{Tăng số lần lặp mỗi ô} (sai số chuẩn giảm theo $\sqrt{n}$) và
   \textbf{mở rộng dải điểm gãy}.
 \item \textbf{Thêm baseline của chính tác giả benchmark} cho bài 2.
 \item \textbf{Gỡ thư mục dữ liệu thô về layout phẳng} và \textbf{phát hành bộ dữ liệu}.
\end{enumerate}

\appendix
\section{Tái lập từng con số}
\begin{footnotesize}
\begin{verbatim}
# sai so diem gay + hai moc nen
python papers/p1_du_phong/experiments/feasibility/assess_predictive_performance.py
# benchmark parser -- can goi LLM that
python papers/p1_du_phong/experiments/parser/parser_benchmark_suite.py
# tuong duong hai ngan xep
python papers/p1_du_phong/experiments/feasibility/doi_chieu_hai_he.py
# phep kiem can thiep cho Tang 2.5
python experiments/chua_phan_loai/rq7_interventional_validity.py
# hop dong du lieu -- phai 0 FAIL
python papers/p1_du_phong/experiments/collect/data_contract_check.py \
    --dir data/raw/SS-PROSP2/SS-PROSP2 --role limits
\end{verbatim}
\end{footnotesize}

\noindent\textbf{Mốc hồi quy.} Script đánh giá P3 chạy \emph{không cờ nào} phải in
đúng \texttt{119.7 / 87.5 / 72.6}. Lệch nghĩa là có gì đã đổi ngoài ý muốn.

\luuy{\textbf{Lưu ý đường dẫn dữ liệu.} Trên máy hiện tại thư mục dữ liệu thô
\emph{lồng thêm một cấp}: thư mục thật là \texttt{data/raw/SS-PROSP2/SS-PROSP2}. Hàm
đọc dữ liệu quét theo mẫu hai cấp, nên lệnh như trong runbook
(\texttt{-{}-dir data/raw/SS-PROSP2}) sẽ tìm thấy \emph{0} lượt đo.}

\end{document}
